\chapter{Discussion, Future Work \& Conclusion}
\label{discussion}

\section{Discussion}

This thesis asked whether post-\gls{sirt} \gls{y90} imaging can be improved by using PET, SPECT and CT together during reconstruction rather than independently, and, if so, how that coupling should be expressed. The answer developed across Chapters~\ref{chapterlabel3}--\ref{chapter:paper} is that it can, provided three things are addressed together: the Bremsstrahlung SPECT forward model, the form of the inter-modality coupling, and the optimisation of the resulting objective. Treating any one of these in isolation was found to be insufficient.

The clearest result of the thesis is that anatomical guidance and functional coupling are complementary rather than interchangeable. In the bootstrapped \gls{nema} study of Chapter~\ref{chapter:paper}, \dtv{} (directional guidance without coupling) and \wtnv{} (coupling without directional guidance) produced broadly similar PET appearances and similar recovery, and neither recovered the smallest spheres at the matched operating point. Recovery improved significantly only when both mechanisms acted together, for all sphere sizes $\leq\SI{22}{\milli\metre}$. This is consistent with the intended mechanism of \wdtnv: the higher-count SPECT image provides functional support for the location of activity transitions, while the $\mu$-derived direction field determines which gradients are penalised at those locations. Neither piece of side information is sufficient on its own in this count regime, which is a useful result for prior design beyond the specific application considered here.

The deconvolution study of Chapter~\ref{chapter:deconvolution_pvc} anticipated this outcome in a simpler setting, and also isolated a second distinction that matters for quantification. \gls{krl} and \gls{hkrl} carry no explicit penalty, so the stopping iteration forms part of the regularisation. The final estimate is determined by the selected iteration, not by minimising an explicit objective. In that study the stopping point had to be selected using the \gls{nrmse} against a known ground truth for the sphere phantom, and by visual assessment for the brain data, neither of which is available in a clinical reconstruction. The consequences are an iteration-dependent bias and, since the iterate retains the influence of its starting point, a marked dependence on the initialising image. \gls{rl}, \gls{krl} and \gls{hkrl} were all sensitive to whether reconstruction was initialised with a uniform image or the blurred \gls{pet} image, to the extent that the best unregularised \gls{rl} result under blurred initialisation occurred at iteration zero. \gls{rldtv}, by contrast, reached the same optima under both initialisations and achieved the lowest \gls{nrmse}, the highest aggregate \gls{rc} and the lowest background \gls{cov} in either case.

This distinction persists in the full reconstruction setting. \gls{shkem} in Chapter~\ref{chapter:paper} is similarly evaluated at a selected subiteration, chosen to approximately match the background \gls{cov} of the vendor reconstruction, whereas \wdtnv{} is iterated towards the minimiser of a fixed objective and its operating point is set by the regularisation weights. Both approaches require a choice to be made, but the penalised formulation states that choice as part of the model, which makes it transferable between datasets and, as Chapter~\ref{chapter:optimisation} demonstrates, separable from the behaviour of the optimiser. This is the argument made in Chapter~\ref{chapter:reconstruction} for preferring \gls{map} estimation to early-stopped \gls{ml}, and the deconvolution study is where it was tested directly. The deconvolution study also distinguishes two uses of anatomical guidance: \gls{krl} and \gls{hkrl} build it into the kernelised image model, whereas \gls{rldtv} uses it in an explicit regularisation term. The same distinction reappears at the small sphere sizes of Chapter~\ref{chapter:paper}, where \wdtnv{} separates from \gls{shkem}.

The treatment of CT follows from the same reasoning. Chapter~\ref{chapter:wdtnv} separated two roles that are easily confounded in multimodality reconstruction, using CT directionally through $\Pi_{m,n}$ and handling modality scale separately through $B_n$. Because CT never enters as a third coupled intensity channel, the voxel-wise Jacobian remains a two-modality Jacobian and the key spectral operation reduces to a $2\times2$ inverse square root, which is what makes the prior and its curvature models affordable at clinical image sizes. The cost of this restraint is that CT contributes orientation only and cannot contribute contrast; the benefit is that the prior cannot transfer CT contrast into the emission images, which is the principal failure mode of intensity-coupled anatomical priors.

Chapter~\ref{chapterlabel3} is best read as a precondition for this coupling rather than as an independent contribution. A joint prior shares whatever each channel contains, including systematic modelling error, so coupling a poorly modelled Bremsstrahlung SPECT channel to a count-limited PET channel risks propagating bias in both directions. Fixing the additive non-primary term through an under-relaxed, bounded reconstruct--simulate refinement, and representing net response broadening by a calibrated effective range kernel, gives a SPECT model that is stable enough to be coupled. The wider point is that synergistic reconstruction raises, rather than lowers, the standard required of each individual forward model.

Optimisation proved to be an intrinsic part of the contribution rather than an implementation detail. In Chapter~\ref{chapter:optimisation}, prior-aware preconditioning was necessary, not merely beneficial, for the count-limited PET channel: data-only \gls{bsrem} scaling diverged under the tested step-size schedule, while all four prior curvature models remained bounded and continued to reduce regional error. SPECT was considerably more robust, and data-only scaling remained competitive there, indicating that the same prior induces very different effective conditioning in the two channels. The four curvature models behaved similarly to one another, so the practically useful conclusion is the binary one: the preconditioner must see the prior. This aligns with independent work on variance-reduced PET reconstruction with prior-aware preconditioning \cite{ehrhardtFastPETReconstruction2025a}. A comparable observation applies to subset organisation, where the schedule of prior application mattered considerably more than how the data subsets were grouped, with only the folded schedule converging at matched average prior weight.

Clinically, \wdtnv{} yielded significantly higher lesion \gls{tbr} than \dtv, \wtnv{} and \gls{shkem} across 45 lesions in nine patients at approximately matched background \gls{cov}. The margin over \gls{shkem} is the smallest of the three, which is expected: \gls{shkem} incorporates anatomical information through CT-informed kernel features and is a strong comparator rather than a naive baseline. In the phantom it matched \wdtnv{} for spheres of \SI{22}{\milli\metre} and above, separating only at \SI{13}{\milli\metre} and \SI{10}{\milli\metre}, which is precisely the regime the prior was designed to address. The reproducibility analysis shows the expected counterpart to this contrast recovery: at the matched operating point, the more strongly smoothing \dtv{} and \wtnv{} baselines give lower voxel-wise \gls{cov} than \wdtnv, and the appropriate reading is that \wdtnv{} spends less of its regularisation budget on suppressing structure that it can instead localise.

Lesion-level responses were heterogeneous, and this is informative rather than simply a limitation. Lesions with clear anatomical correlates, such as adjacent calcification or a liver capsule boundary, benefited most, while lesions without a corresponding structure in the $\mu$-map showed limited improvement. This is the designed behaviour: CTAC guidance steers orientation preferences where anatomy is present and defers to the functional data where it is not, and the prior cannot invent edges that neither modality supports. It also suggests a practical route to case selection, since the directional weighting field $\|\xi\|$ is available independently of the reconstruction method and could in principle indicate, once lesions have been identified on a routine clinical reconstruction, which of them are likely to benefit.

Placing these results in the wider clinical context is important for interpreting their significance. The case for joint reconstruction does not rest on either modality being unusable: PET- and SPECT-based mean absorbed doses are highly correlated across liver \glspl{voi} \cite{kappadathQuantitativeEvaluation90YPET2024}. What that comparison also reports, however, is a systematic modality-dependent offset, with PET yielding on average $14\%$ higher tumour mean dose and lower perfused-normal-liver mean dose than SPECT. A discrepancy of that form, concentrated at the lesion level and in the direction expected from the differing resolution and partial-volume behaviour of the two modalities, is precisely the quantity that reconstruction can influence, and it is where the improvements reported here are concentrated. As dosimetry-guided prescription for \gls{sirt} becomes more formalised \cite{salemClinicalDosimetricConsiderations2019} and dosimetry itself more widely implemented \cite{gearEANMEnablingGuide2023}, the accuracy of lesion-level activity estimation matters more, not less. Continuing improvements in scanner sensitivity will reduce the severity of the count-limited PET regime at some centres, but they do not address the modelling difficulties specific to Bremsstrahlung SPECT, and most centres performing post-\gls{sirt} imaging will continue to rely on both modalities.

\subsection{Limitations}
\label{subsec:discussion_limitations}

The most practically significant limitation is the sensitivity of the directional field to CTAC quality. Because $\xi$ is derived from $D\mu$, coherent CT artefacts such as arms-down streaking are reproduced in $\|\xi\|$ and provide a pathway for artefact transfer into the emission reconstructions, as observed in the training patient of Chapter~\ref{chapter:paper}. The stabiliser $\eta$ limits amplification of small gradients but cannot suppress structured, high-contrast artefacts. This is a property of deriving guidance from the raw attenuation map rather than of the prior itself, and it is addressable; where a low-dose CTAC is used for guidance, acquisition with arms raised where feasible is preferable. Relatedly, the sensitivity of \wdtnv{} to residual PET/SPECT/CT registration error was not characterised. A prior that rewards aligned edges is, by construction, the class of method most exposed to misregistration, so this deserves explicit study.

The evaluation protocol also carries limitations that follow from the decision to preserve direct comparability with \gls{shkem}. Regularisation weights were selected once on a single training patient and then fixed, which avoids repeated tuning on the evaluation cohort but does not establish the robustness of those weights across acquisition protocols and lesion presentations. Quantitative endpoints were restricted to PET, so the jointly reconstructed SPECT images, although visually less noisy while retaining lesion conspicuity, were assessed qualitatively only. The endpoints themselves are image-quality surrogates: recovery, \gls{tbr} and \gls{cov} are the quantities that reconstruction most directly influences, but the clinical motivation is absorbed dose, and no dose-level or outcome-level analysis was performed. Finally, the clinical cohort is single-centre and retrospective, and the thesis compares against variational and kernel-based methods but not against learned reconstruction or learned correction approaches \cite{hashimotoPETImageReconstruction2022,jia90YSPECTScatter2023}, which are the most likely external competitors for this problem.

\section{Future Work}
\label{sec:future_work}

\subsection{Relative Nuclear Variation}
\label{subsec:future_rnv}

The modality weights used by \wdtnv{} compensate for the markedly different global scales of the PET and SPECT images, but do not adapt to spatial variation in activity within either image. One possible extension, inspired by the \acrfull{rdp} \cite{nuytsConcavePriorPenalizing2002}, would scale the coupled edge magnitude by a local activity measure. Let $r_n(\bm{z})=\|\tilde{\mathcal{G}}_n(\bm{z})\|_*$ denote the nuclear magnitude of the weighted directional Jacobian used by \wdtnv{}, and let $s_n(\bm{z})\geq0$ denote a weighted local activity sum over the same finite-difference neighbourhood. A proposed \acrfull{rnv} penalty is
\begin{equation}
    \mathcal{R}_{\mathrm{RNV}}(\bm{z})
    :=
    V_{\mathrm{vox}}
    \sum_{n=1}^{N}
    \frac{r_n(\bm{z})^2}
    {s_n(\bm{z})+\gamma_{\mathrm{RNV}}r_n(\bm{z})
        +\varepsilon_{\mathrm{RNV}}},
    \qquad
    \gamma_{\mathrm{RNV}}>0,\quad
    \varepsilon_{\mathrm{RNV}}>0.
    \label{eq:rnv_prior}
\end{equation}
The construction is ``relative'' because the penalty on a coupled edge is moderated by the local activity scale. For fixed non-negative modality weights on the non-negative image domain, the scalar density $r^2/(s+\gamma_{\mathrm{RNV}}r+\varepsilon_{\mathrm{RNV}})$ is jointly convex in $(r,s)$ and non-decreasing in $r$; since $r_n$ is convex and $s_n$ is affine, the proposed penalty is convex under these assumptions. This construction was not implemented or evaluated in this thesis, and whether such local scaling improves the bias--noise trade-off without weakening cross-modality regularisation inappropriately remains to be established.

\subsection{Robustness and Clinical Translation}
\label{subsec:future_translation}

Three methodological directions follow directly from the limitations above. First, the directional field should be made robust to CTAC quality. Estimating $\xi$ from a denoised, artefact-corrected or segmentation-derived attenuation map, rather than from $D\mu$ alone, would retain the anatomical orientation information while removing the pathway by which coherent streak artefacts enter the emission reconstruction. Second, the sensitivity of \wdtnv{} to residual registration error should be quantified directly, for example by applying controlled rigid and deformable perturbations to the SPECT-to-PET warp and measuring the resulting lesion-level bias; a joint reconstruction--registration formulation is the natural extension if that sensitivity proves material. Third, the optimisation results of Chapter~\ref{chapter:optimisation} indicate that a single shared step-size schedule is not equally well calibrated to both modalities, since increasing the initial step improved SPECT convergence while destabilising the count-limited PET reconstruction. Modality-dependent initial steps or relaxation rates are therefore a straightforward and likely worthwhile refinement.

Beyond this, two extensions would broaden the applicability of the framework. The present study addresses the favourable case in which PET and SPECT image the same underlying \gls{y90} microsphere distribution. Settings with shared anatomy but non-identical functional distributions, such as \gls{maa} planning versus post-treatment imaging, multi-nuclide studies, or multiple timepoints, would require the coupling strength to be relaxed or made spatially adaptive so that modality-specific or temporally evolving uptake is not suppressed. Finally, the endpoints should be advanced from image quality to absorbed dose. Propagating the reconstructions through voxel-level dose kernels and evaluating dose--volume endpoints, ideally in a multi-centre cohort with response data, would test whether the lesion-level improvements reported here translate into more reliable dosimetry, which is the ultimate motivation for the work.

Two further comparisons would strengthen the case. The effective range kernel of Chapter~\ref{chapterlabel3} is isotropic and water-equivalent, so replacing it with an anisotropic, material-dependent Bremsstrahlung model such as that of Lim \textit{et al.} \cite{lim_y-90_2018} would test how much of the residual SPECT modelling error remains attributable to that approximation. Separately, the comparators used throughout this thesis are variational and kernel-based, and a comparison against learned reconstruction and learned correction methods \cite{hashimotoPETImageReconstruction2022,jia90YSPECTScatter2023} would establish where an interpretable, physically constrained prior stands relative to the data-driven alternatives now being developed for this problem.

\section{Conclusion}
\label{sec:thesis_conclusion}

This thesis set out to improve quantitative post-\gls{sirt} \gls{y90} imaging by using PET, SPECT and CT together during reconstruction rather than independently, and to do so in a form that remains tractable at clinical image sizes. Four contributions support that aim.

First, Chapter~\ref{chapterlabel3} established the Bremsstrahlung SPECT model components required before the two emission modalities can safely be coupled. A domain-shifted learned projection estimate was used to initialise object- and acquisition-specific SIMIND simulation of non-primary counts; an under-relaxed update remained bounded over the refinement budget where full replacement did not; and a calibrated effective Gaussian range kernel brought projected response widths substantially closer to GATE and SIMIND across three energy windows. The resulting model was frozen and carried forward unchanged.

Second, Chapter~\ref{chapter:wdtnv} introduced \wdtnv, a smoothed, weighted and directional form of \gls{tnv}. Charbonnier smoothing of the singular values yields a differentiable spectral functional whose Jacobian-space gradient can be written in Gram-matrix form, reducing the key spectral operation for two coupled emission modalities to a $2\times2$ inverse square root and removing the need for a per-voxel \gls{svd}. The formulation separates anatomical direction, applied through $\Pi_{m,n}$, from modality scale, applied through $B_n$, so that CT steers the penalty without entering as a coupled intensity channel. The design principle behind that choice was established beforehand in the controlled deconvolution setting of Chapter~\ref{chapter:deconvolution_pvc}, where an explicit directional penalty outperformed kernelised image models and, unlike them, gave solutions that depended on neither the initialisation nor a ground-truth-informed stopping rule.

Third, Chapter~\ref{chapter:optimisation} treated the optimisation of the coupled objective as part of the contribution rather than as an implementation concern. The PET bed positions are reconstructed jointly, so that the edge-preserving prior acts consistently across their overlap instead of being applied independently before fusion. With the average prior weight per data pass matched, the prior-application schedule mattered more than the subset organisation, and only the folded schedule converged. Prior-aware preconditioning was required for stable convergence of the count-limited PET channel, where data-only scaling diverged, although the four prior curvature models examined behaved similarly to one another.

Fourth, Chapter~\ref{chapter:paper} evaluated the complete framework on \gls{y90} data. On a bootstrapped \gls{nema} phantom, \wdtnv{} significantly improved recovery of spheres of \SI{22}{\milli\metre} and below relative to \dtv{} and \wtnv, and improved on \gls{shkem} at the two smallest sphere sizes, at a matched noise operating point. In a cohort of 45 lesions across nine post-\gls{sirt} patients, \wdtnv{} gave significantly higher lesion \gls{tbr} than all three comparators at approximately matched background \gls{cov}. The simulation and reconstruction software developed to support this work is documented in Appendix~\ref{app:software} and released so that the workflows can be reproduced.

The position these results support is a specific one. Anatomical information is most safely used to steer the geometry of a penalty rather than to constrain image intensity or to restrict the space in which the solution may lie; functional coupling and anatomical guidance are complementary, and in this count regime neither is sufficient alone; and a prior strong enough to deliver that benefit must be accompanied by an optimiser that accounts for its curvature. The clinical case is not yet complete, since the evaluation reported here is single-centre, rests on image-quality endpoints rather than absorbed dose, and does not establish robustness to CT artefact or residual registration error. These are the natural next steps set out in Section~\ref{sec:future_work} rather than obstacles in principle, and the framework developed here provides a practical basis on which to pursue them.
